\chapter{System Design and Implementation}

\section{System Architecture}
The ShifaMart+ AI Agent follows a modular architecture with clear separation of concerns. The system consists of several interconnected components that work together to provide comprehensive healthcare assistance.

\subsection{High-Level Architecture}
The system architecture is divided into the following layers:

\begin{enumerate}
    \item \textbf{Presentation Layer}: Web interface and API endpoints
    \item \textbf{Application Layer}: Business logic and orchestration
    \item \textbf{AI/ML Layer}: Machine learning models and NLP processors
    \item \textbf{Data Layer}: Data preprocessing and model storage
\end{enumerate}

\subsection{Component Diagram}
The main components of the system include:

\begin{itemize}
    \item \textbf{Disease Prediction Model}: Ensemble of XGBoost and Random Forest classifiers
    \item \textbf{NLP Processor}: Natural language understanding for symptom extraction
    \item \textbf{Severity Detection Model}: Gradient Boosting classifier for severity assessment
    \item \textbf{First Aid System}: Rule-based first-aid guidance system
    \item \textbf{Conversational Agent}: State machine-based chat interface
    \item \textbf{Specialist Mapper}: Disease-to-specialist recommendation system
    \item \textbf{Maps Integration}: Location-based doctor search
    \item \textbf{API Server}: FastAPI-based RESTful API
\end{itemize}

\section{Data Preprocessing}
The data preprocessing module handles loading and transformation of medical datasets.

\subsection{Dataset Structure}
The system uses four main datasets:

\begin{enumerate}
    \item \textbf{Main Dataset} (\texttt{dataset.csv}): Contains 4,922 symptom-disease mappings
    \item \textbf{Symptom Description} (\texttt{symptom\_Description.csv}): Disease descriptions
    \item \textbf{Symptom Precaution} (\texttt{symptom\_precaution.csv}): Precautions for each disease
    \item \textbf{Symptom Severity} (\texttt{Symptom-severity.csv}): Severity weights for symptoms
\end{enumerate}

\subsection{Preprocessing Steps}
The preprocessing pipeline includes:

\begin{enumerate}
    \item \textbf{Data Loading}: Reading CSV files and handling missing values
    \item \textbf{Symptom Normalization}: Converting symptom names to standardized format
    \item \textbf{Feature Encoding}: Creating binary feature vectors from symptoms
    \item \textbf{Label Encoding}: Encoding disease names to numeric labels
    \item \textbf{Data Splitting}: Training/validation/test set division
\end{enumerate}

\section{Disease Prediction Model}
The disease prediction model uses an ensemble approach combining XGBoost and Random Forest classifiers.

\subsection{Model Architecture}
The ensemble model consists of:

\begin{itemize}
    \item \textbf{Random Forest Classifier}:
    \begin{itemize}
        \item 100 estimators
        \item Maximum depth: 15
        \item Minimum samples split: 5
        \item Class weight: balanced
    \end{itemize}
    
    \item \textbf{XGBoost Classifier}:
    \begin{itemize}
        \item 100 estimators
        \item Maximum depth: 6
        \item Learning rate: 0.1
        \item Objective: multi:softprob
        \item Tree method: hist (memory efficient)
    \end{itemize}
    
    \item \textbf{Voting Classifier}: Soft voting using probability distributions
\end{itemize}

\subsection{Training Process}
The training process includes:

\begin{enumerate}
    \item Data preprocessing and feature extraction
    \item Train-test split (80-20) with stratification
    \item Individual model training
    \item Ensemble model creation and training
    \item Model evaluation using accuracy and classification report
    \item Model persistence using joblib
\end{enumerate}

\subsection{Rule-Based Enhancements}
To improve prediction accuracy, the system includes rule-based corrections:

\begin{itemize}
    \item \textbf{Symptom-Disease Pattern Boosting}: Increases probability for common symptom-disease combinations
    \item \textbf{Serious Disease Filtering}: Reduces probability of serious diseases when only generic symptoms are present
    \item \textbf{Duration-Based Adjustments}: Considers symptom duration for better predictions
\end{itemize}

\section{Natural Language Processing}
The NLP processor extracts symptoms from natural language descriptions.

\subsection{Processing Pipeline}
The NLP pipeline includes:

\begin{enumerate}
    \item \textbf{Text Preprocessing}:
    \begin{itemize}
        \item Lowercase conversion
        \item Contraction expansion
        \item Punctuation normalization
        \item Whitespace normalization
    \end{itemize}
    
    \item \textbf{Symptom Extraction}:
    \begin{itemize}
        \item Synonym matching (132+ symptom synonyms)
        \item Direct symptom matching
        \item Fuzzy matching for common symptom words
    \end{itemize}
    
    \item \textbf{Metadata Extraction}:
    \begin{itemize}
        \item Duration extraction (days, weeks, hours, months)
        \item Severity modifier detection (mild, moderate, severe)
        \item Negation detection
    \end{itemize}
\end{enumerate}

\subsection{Synonym Mapping}
The system maintains a comprehensive synonym dictionary mapping natural language phrases to standardized symptom names. For example:

\begin{itemize}
    \item "fever", "high temperature", "pyrexia" $\rightarrow$ \texttt{high\_fever}
    \item "stomach pain", "tummy ache", "abdominal pain" $\rightarrow$ \texttt{stomach\_pain}
    \item "shortness of breath", "difficulty breathing" $\rightarrow$ \texttt{breathlessness}
\end{itemize}

\section{Severity Detection Model}
The severity detection model classifies symptoms into four severity levels.

\subsection{Severity Levels}
\begin{enumerate}
    \item \textbf{MILD}: Low-risk symptoms requiring home care
    \item \textbf{MODERATE}: Symptoms that should be monitored
    \item \textbf{SEVERE}: Symptoms requiring medical attention
    \item \textbf{EMERGENCY}: Critical symptoms requiring immediate care
\end{enumerate}

\subsection{Model Approach}
The severity model uses a hybrid approach:

\begin{itemize}
    \item \textbf{Rule-Based Component}:
    \begin{itemize}
        \item Symptom weight scoring
        \item Dangerous combination detection
        \item Duration-based severity adjustment
    \end{itemize}
    
    \item \textbf{ML Component}:
    \begin{itemize}
        \item Gradient Boosting Classifier
        \item Synthetic training data generation
        \item Confidence-based prediction
    \end{itemize}
\end{itemize}

\subsection{Emergency Detection}
The system identifies emergencies through:

\begin{itemize}
    \item Critical symptom detection (chest pain, breathlessness, coma, etc.)
    \item Dangerous symptom combinations (e.g., chest pain + breathlessness = heart attack)
    \item High severity scores from weighted symptom analysis
\end{itemize}

\section{First Aid System}
The first aid system provides step-by-step guidance for medical emergencies.

\subsection{Emergency Types Covered}
The system covers 12+ emergency types including:

\begin{itemize}
    \item Heart attack
    \item Stroke
    \item Breathing difficulty
    \item Choking
    \item Severe bleeding
    \item Burns
    \item Seizures
    \item Unconsciousness
    \item Anaphylaxis
    \item Poisoning
    \item High fever emergency
    \item Severe dehydration
\end{itemize}

\subsection{First Aid Guide Structure}
Each first aid guide includes:

\begin{itemize}
    \item Title and description
    \item Emergency contact numbers
    \item Step-by-step instructions
    \item Warnings and contraindications
    \item "Do NOT" guidelines
    \item When to call emergency services
    \item Additional notes
\end{itemize}

\section{Conversational AI Agent}
The conversational agent provides an intuitive interface for symptom collection and analysis.

\subsection{State Machine Design}
The agent uses a state machine with the following states:

\begin{enumerate}
    \item \textbf{GREETING}: Initial welcome and instructions
    \item \textbf{COLLECTING\_SYMPTOMS}: Active symptom collection
    \item \textbf{ASKING\_DURATION}: Duration inquiry
    \item \textbf{CONFIRMING\_SYMPTOMS}: Symptom confirmation
    \item \textbf{SHOWING\_RESULTS}: Displaying predictions and recommendations
    \item \textbf{EMERGENCY\_DETECTED}: Emergency response mode
    \item \textbf{ASKING\_CITY}: Location input for doctor search
    \item \textbf{ENDED}: Session termination
\end{enumerate}

\subsection{Session Management}
Each user session maintains:

\begin{itemize}
    \item Collected symptoms
    \item Symptom duration
    \item Disease predictions
    \item Severity assessment
    \item Conversation history
    \item User location (if provided)
\end{itemize}

\section{Specialist Recommendation}
The system recommends appropriate medical specialists based on predicted diseases.

\subsection{Specialist Mapping}
The system maps diseases to specialists including:

\begin{itemize}
    \item Cardiologist (heart conditions)
    \item Neurologist (neurological conditions)
    \item Dermatologist (skin conditions)
    \item Gastroenterologist (digestive conditions)
    \item Pulmonologist (respiratory conditions)
    \item And 10+ other specialties
\end{itemize}

\subsection{Recommendation Logic}
Specialist recommendation uses:

\begin{enumerate}
    \item Primary: Disease-based mapping (most accurate)
    \item Fallback: Symptom-based mapping
    \item Integration with location services for finding nearby doctors
\end{enumerate}

\section{API Implementation}
The system is implemented as a RESTful API using FastAPI.

\subsection{API Endpoints}
Key endpoints include:

\begin{itemize}
    \item \texttt{POST /predict/symptoms}: Predict from symptom list
    \item \texttt{POST /predict/text}: Predict from natural language text
    \item \texttt{POST /chat}: Conversational agent interaction
    \item \texttt{POST /severity/check}: Severity assessment
    \item \texttt{POST /first-aid/for-symptoms}: First aid guidance
    \item \texttt{GET /symptoms}: List all symptoms
    \item \texttt{GET /diseases}: List all diseases
    \item \texttt{GET /health}: System health check
\end{itemize}

\subsection{Response Format}
API responses follow a consistent structure:

\begin{lstlisting}[language=Python, caption=Example API Response]
{
    "predictions": [
        {
            "disease": "Common Cold",
            "probability": 0.85,
            "confidence_percent": "85.0%",
            "description": "Viral infection...",
            "precautions": ["rest", "hydration"],
            "severity_score": 2.5,
            "severity_level": "MILD"
        }
    ],
    "severity": {
        "level": "MILD",
        "score": 2.5,
        "confidence": 0.8,
        "reason": "Mild symptoms...",
        "is_emergency": false
    },
    "recommended_specialist": {
        "name": "General Practitioner",
        "description": "...",
        "icon": "👨‍⚕️"
    }
}
\end{lstlisting}

\section{Implementation Technologies}
The system is implemented using:

\begin{itemize}
    \item \textbf{Python 3.8+}: Primary programming language
    \item \textbf{FastAPI}: Web framework for API
    \item \textbf{scikit-learn}: Machine learning library
    \item \textbf{XGBoost}: Gradient boosting framework
    \item \textbf{pandas/numpy}: Data manipulation
    \item \textbf{joblib}: Model serialization
    \item \textbf{uvicorn}: ASGI server
\end{itemize}

\section{Model Training and Evaluation}
\subsection{Training Process}
Models are trained using:

\begin{itemize}
    \item 80-20 train-test split
    \item Stratified sampling for balanced classes
    \item Cross-validation for hyperparameter tuning
    \item Model persistence for production use
\end{itemize}

\subsection{Evaluation Metrics}
The system is evaluated using:

\begin{itemize}
    \item Accuracy: Overall prediction correctness
    \item Precision: Correct positive predictions
    \item Recall: Ability to find all positive cases
    \item F1-Score: Harmonic mean of precision and recall
    \item Classification Report: Per-disease metrics
\end{itemize}

\section{System Deployment}
The system can be deployed as:

\begin{itemize}
    \item Standalone API server
    \item Docker container
    \item Cloud service (AWS, Azure, GCP)
    \item Integration with web/mobile applications
\end{itemize}
